\section{Wedderburn decomposition of the fixed-point algebra (continued)}

\begin{theorem}[Isotypic-component decomposition]%
    \label{thm:starSubalgebra_orthogonal_isotypic_decomposition}
    \lean{StarSubalgebra.exists_orthogonal_isotypic_decomposition}
    \leanok
    \uses{thm:starSubalgebra_orthogonal_irreducible_decomposition,
        lem:starSubalgebra_isOrtho_sSup_not_sameIsotype,
        lem:starSubalgebra_sameIsotype_of_le_sSup,
        lem:starSubalgebra_sameIsotype_refl, lem:starSubalgebra_sameIsotype_symm,
        lem:starSubalgebra_sameIsotype_trans}
    Let $S$ be a $*$-subalgebra of $\MN{D}$. Then $\C^{D}$ is the supremum of a finite family
    $\mathcal{C}$ of pairwise orthogonal subspaces, the isotypic components, and each component
    $C \in \mathcal{C}$ is the supremum of a finite, nonempty class $\mathcal{D}_C$ of pairwise
    orthogonal subspaces irreducible under $S$ that are pairwise of the same type:
    \[
        C = \bigvee_{W \in \mathcal{D}_C} W.
    \]
    Moreover, for distinct components $C \neq C'$, no subspace irreducible under $S$ contained
    in $C$ is of the same type as any subspace irreducible under $S$ contained in $C'$.
\end{theorem}

\begin{proof}
    \leanok
    \uses{thm:starSubalgebra_orthogonal_irreducible_decomposition,
        lem:starSubalgebra_isOrtho_sSup_not_sameIsotype,
        lem:starSubalgebra_sameIsotype_of_le_sSup,
        lem:starSubalgebra_sameIsotype_refl, lem:starSubalgebra_sameIsotype_symm,
        lem:starSubalgebra_sameIsotype_trans}
    The orthogonal irreducible decomposition
    (Theorem~\ref{thm:starSubalgebra_orthogonal_irreducible_decomposition}) supplies a finite,
    pairwise orthogonal family $\mathcal{D}$ of irreducible pieces with
    $\bigvee_{W \in \mathcal{D}} W = \C^{D}$. For each piece $W$ let $\mathcal{D}_W$ be the class
    of pieces of $\mathcal{D}$ of the same type as $W$, and let $C_W$ be its supremum. Each piece
    lies in its own class by reflexivity
    (Lemma~\ref{lem:starSubalgebra_sameIsotype_refl}), so each class is nonempty, the classes
    cover $\mathcal{D}$, and the components span $\C^{D}$. Each class is contained in
    $\mathcal{D}$, so its pieces are pairwise orthogonal. Two pieces of one class are of the
    same type with each other through the representative, by symmetry and transitivity
    (Lemmas~\ref{lem:starSubalgebra_sameIsotype_symm}
    and~\ref{lem:starSubalgebra_sameIsotype_trans}). If two representatives are of the same type,
    transitivity makes their classes equal, hence their components equal; so two distinct
    components come from representatives $W, W'$ that are not of the same type. Write $\sim$ for
    the same-type relation. For $a \in \mathcal{D}_W$ and $b \in \mathcal{D}_{W'}$ one has
    $W \sim a$ and $W' \sim b$ by definition of the classes; were $a \sim b$, symmetry and
    transitivity would give
    \[
        W \sim a \sim b \sim W',
    \]
    hence $W \sim W'$, contradicting that $W$ and $W'$ are not of the same type. So no piece of
    $\mathcal{D}_W$ is of the same type as any piece of $\mathcal{D}_{W'}$, and the components are
    orthogonal (Lemma~\ref{lem:starSubalgebra_isOrtho_sSup_not_sameIsotype}).

    Finally, let $V \le C_W$ and $V' \le C_{W'}$ be irreducible under $S$ with
    $C_W \neq C_{W'}$. By Lemma~\ref{lem:starSubalgebra_sameIsotype_of_le_sSup} there are
    pieces $a \in \mathcal{D}_W$ and $b \in \mathcal{D}_{W'}$ with $V \sim a$ and $V' \sim b$.
    Were $V \sim V'$, symmetry and transitivity would give
    \[
        W \sim a \sim V \sim V' \sim b \sim W',
    \]
    contradicting again that $W$ and $W'$ are not of the same type.
\end{proof}

The pieces of one class are unitarily identified copies of a single irreducible piece. An
orthonormal basis of one piece therefore spreads, through the inner-product-preserving
intertwiners, to an adapted orthonormal basis of the whole component, and in this basis every
member of $S$ acts by one matrix on the irreducible index, identically across the multiplicity
index. This is the tensor factorization of the isotypic components behind the block form
$\MN{d_k} \otimes \Id_{m_k}$ of~\cite[Theorem~6.14]{Wolf2012Quantum}.

\begin{theorem}[Adapted orthonormal basis of an isotypic component]%
    \label{thm:starSubalgebra_adapted_orthonormal_basis}
    \lean{StarSubalgebra.exists_adapted_orthonormal_family_of_sameIsotype}
    \leanok
    \uses{def:starSubalgebra_irreducibleSubspace, def:starSubalgebra_sameIsotype}
    Let $S$ be a $*$-subalgebra of $\MN{D}$ and let $\mathcal{D}$ be a finite, nonempty family
    of pairwise orthogonal subspaces of $\C^{D}$, irreducible under $S$ and pairwise of the
    same type. Write $m$ for the number of pieces of $\mathcal{D}$; all pieces share one
    dimension $d$. Then the component $C = \bigvee_{W \in \mathcal{D}} W$ has an orthonormal
    basis $(f_{i,j})_{0 \le i < m,\, 0 \le j < d}$ such that for every $A \in S$ there is a
    matrix $B \in \MN{d}$ with, for all $i < m$ and $j < d$,
    \[
        A\,f_{i,j} \;=\; \sum_{j'=0}^{d-1} B_{j'j}\, f_{i,j'}.
    \]
    In the adapted basis the member $A$ acts on $C$ by the matrix $B$ on the irreducible
    index, identically across the multiplicity index; after reordering the indices this is the
    block $\MN{d} \otimes \Id_{m}$ of~\cite[Theorem~6.14]{Wolf2012Quantum}. The vectors
    $(f_{i,j})_{0 \le j < d}$ of the $i$-th copy span one of the pieces of $\mathcal{D}$.
\end{theorem}

\begin{proof}
    \leanok
    \uses{lem:starSubalgebra_sameIsotype_isometry, lem:starSubalgebra_sameIsotype_refl}
    Fix a piece $W \in \mathcal{D}$, of dimension $d$, with an orthonormal basis
    $e_0, \ldots, e_{d-1}$, and enumerate $\mathcal{D} = \{W'_0, \ldots, W'_{m-1}\}$. Every
    piece of $\mathcal{D}$ is of the same type as $W$, by hypothesis for the other pieces and
    by reflexivity (Lemma~\ref{lem:starSubalgebra_sameIsotype_refl}) for $W$ itself, so
    Lemma~\ref{lem:starSubalgebra_sameIsotype_isometry} provides intertwiners $u_i$ mapping
    $W$ onto $W'_i$ and preserving the inner product on $W$. Set $f_{i,j} = u_i(e_j)$.
    Within one copy the vectors $f_{i,0}, \ldots, f_{i,d-1}$ are orthonormal because $u_i$
    preserves inner products; across copies they are orthogonal because distinct pieces of
    $\mathcal{D}$ are. Since $u_i$ maps $W$ onto $W'_i$, the vectors of the $i$-th copy span
    $W'_i$; in particular each $W'_i$ has dimension $d$, and all the vectors together span $C$.
    For $A \in S$ the piece $W$ is invariant, so $A e_j = \sum_{j'} B_{j'j}\, e_{j'}$ with
    $B_{j'j} = \langle e_{j'}, A e_j \rangle$, and the intertwining identity gives
    \[
        A f_{i,j} = A\,u_i(e_j) = u_i(A e_j)
        = \sum_{j'} B_{j'j}\, u_i(e_{j'}) = \sum_{j'} B_{j'j}\, f_{i,j'},
    \]
    the same matrix $B$ for every copy $i$.
\end{proof}

\begin{theorem}[Spatial tensor factorization of the isotypic components]%
    \label{thm:starSubalgebra_isotypic_tensor_factorization}
    \lean{StarSubalgebra.exists_isotypic_tensor_factorization}
    \leanok
    \uses{thm:starSubalgebra_orthogonal_isotypic_decomposition,
        thm:starSubalgebra_adapted_orthonormal_basis}
    Let $S$ be a $*$-subalgebra of $\MN{D}$. Then $\C^{D}$ is the supremum of a finite family
    of pairwise orthogonal isotypic components, and each component carries an orthonormal
    basis $(f_{i,j})$, indexed by a multiplicity index $i < m$ and an irreducible index
    $j < d$, such that every $A \in S$ acts by
    \[
        A\,f_{i,j} = \sum_{j'=0}^{d-1} B_{j'j}\, f_{i,j'}
    \]
    for a matrix $B \in \MN{d}$ depending only on $A$ and the component. In the adapted bases
    the action of $S$ on each component is by matrix-times-identity blocks, the form
    $\MN{d_k} \otimes \Id_{m_k}$ of~\cite[Theorem~6.14]{Wolf2012Quantum}.
\end{theorem}

\begin{proof}
    \leanok
    \uses{thm:starSubalgebra_orthogonal_isotypic_decomposition,
        thm:starSubalgebra_adapted_orthonormal_basis}
    The isotypic-component decomposition
    (Theorem~\ref{thm:starSubalgebra_orthogonal_isotypic_decomposition}) writes $\C^{D}$ as
    the supremum of finitely many pairwise orthogonal components, each the supremum of a
    finite, nonempty class of pairwise orthogonal irreducible pieces of one type. Applying
    Theorem~\ref{thm:starSubalgebra_adapted_orthonormal_basis} to each class equips each
    component with its adapted orthonormal basis.
\end{proof}

The adapted orthonormal bases of the pairwise orthogonal isotypic components concatenate
into a single orthonormal basis of $\C^{D}$, indexed by a component index $k$, a
multiplicity index $i$, and an irreducible index $j$. In this basis every member of $S$
acts block-diagonally, one matrix-times-identity block per component. A $*$-subalgebra
contains the identity matrix, so every component carries irreducible pieces of $S$: this
is the unital case of the block representation invoked
by~\cite[Theorem~6.14]{Wolf2012Quantum}, without the zero block.

\begin{theorem}[Adapted orthonormal basis of the whole space]%
    \label{thm:starSubalgebra_global_adapted_orthonormal_basis}
    \lean{StarSubalgebra.exists_adapted_orthonormalBasis}
    \leanok
    \uses{def:starSubalgebra_irreducibleSubspace, def:starSubalgebra_sameIsotype}
    Let $S$ be a $*$-subalgebra of $\MN{D}$. There are $K \in \N$, positive dimensions
    $d_0, \ldots, d_{K-1}$ and multiplicities $m_0, \ldots, m_{K-1}$, and an orthonormal
    basis $(f_{k,i,j})_{k < K,\, i < m_k,\, j < d_k}$ of $\C^{D}$ such that for every
    $A \in S$ there are matrices $B_k \in \MN{d_k}$ with, for all $k < K$, $i < m_k$, and
    $j < d_k$,
    \[
        A\,f_{k,i,j} \;=\; \sum_{j'=0}^{d_k-1} (B_k)_{j'j}\, f_{k,i,j'}.
    \]
    In this basis the member $A$ acts on the $k$-th component as
    $\Id_{m_k} \otimes B_k$. Each row $(f_{k,i,j})_{0 \le j < d_k}$ spans a subspace
    irreducible under $S$, and rows drawn from distinct components span subspaces that are
    never of the same type. Only the containment direction is asserted: every member of
    $S$ takes this form. The equality with the block algebra, including the reverse
    inclusion, is Theorem~\ref{thm:starSubalgebra_unitary_block_diagonal_iff}.
\end{theorem}

\begin{proof}
    \leanok
    \uses{thm:starSubalgebra_orthogonal_isotypic_decomposition,
        thm:starSubalgebra_adapted_orthonormal_basis}
    The isotypic-component decomposition
    (Theorem~\ref{thm:starSubalgebra_orthogonal_isotypic_decomposition}) writes $\C^{D}$ as
    the supremum of finitely many pairwise orthogonal isotypic components, each the
    supremum of a finite, nonempty same-type class of pairwise orthogonal irreducible
    pieces. Theorem~\ref{thm:starSubalgebra_adapted_orthonormal_basis} equips the $k$-th
    component with an adapted orthonormal family $(f_{k,i,j})_{i < m_k,\, j < d_k}$
    spanning it, where $m_k$ is the number of pieces of the class, positive because the
    class is nonempty, and $d_k$ is their common dimension, positive because irreducible
    pieces are nonzero. Concatenating the families over $k$ yields an orthonormal family:
    two vectors of one component are orthonormal by the adapted family, and two vectors of
    distinct components are orthogonal because the components are. The concatenated family
    spans $\C^{D}$ because each component is spanned by its family and the components have
    supremum $\C^{D}$; a spanning orthonormal family is an orthonormal basis. The action
    property holds componentwise by the adapted families. Each row of the concatenated
    family spans one of the irreducible pieces of its component's class
    (Theorem~\ref{thm:starSubalgebra_adapted_orthonormal_basis}), and irreducible subspaces
    of distinct components are never of the same type
    (Theorem~\ref{thm:starSubalgebra_orthogonal_isotypic_decomposition}), so neither are the
    row spans of distinct components.
\end{proof}

The change of basis to an orthonormal basis indexed by such triples is implemented by a
unitary matrix, and it carries every matrix acting blockwise on the basis to a
block-diagonal matrix.

\begin{lemma}[Unitary change of basis to block-diagonal form]%
    \label{lem:starSubalgebra_unitary_change_of_basis_block}
    \lean{OrthonormalBasis.exists_unitary_conj_blockDiagonal}
    \leanok
    Let $(f_{k,i,j})_{k < K,\, i < m_k,\, j < d_k}$ be an orthonormal basis of $\C^{D}$.
    There are an identification of the index set of triples $(k, i, j)$ with
    $\{0, \ldots, D-1\}$, realizing $\sum_k d_k m_k = D$, and a unitary $U \in \MN{D}$,
    whose columns are the basis vectors, such that every matrix $A \in \MN{D}$ acting on
    the basis by
    \[
        A\,f_{k,i,j} \;=\; \sum_{j'=0}^{d_k-1} (B_k)_{j'j}\, f_{k,i,j'}
    \]
    for matrices $B_k \in \MN{d_k}$ satisfies
    \[
        U^{\dagger} A\, U \;=\; \bigoplus_{k=0}^{K-1} \Id_{m_k} \otimes B_k.
    \]
\end{lemma}

\begin{proof}
    \leanok
    Let $U$ be the change-of-basis matrix from the standard orthonormal basis of $\C^{D}$
    to $(f_{k,i,j})$, so the columns of $U$ are the basis vectors; a change of basis
    between orthonormal bases is unitary. Both bases have $D$ elements, which identifies
    the index set of triples $(k, i, j)$ with $\{0, \ldots, D-1\}$ and gives
    $\sum_k d_k m_k = D$. For $A$ acting as displayed, the matrix
    $U^{\dagger} A\, U = U^{-1} A\, U$ is the matrix of the operator $x \mapsto Ax$ in the
    basis $(f_{k,i,j})$, with entries
    \[
        \langle f_{k',i',j'},\, A f_{k,i,j} \rangle
        \;=\; \sum_{j''} (B_k)_{j''j} \langle f_{k',i',j'},\, f_{k,i,j''} \rangle
        \;=\; \delta_{k'k}\, \delta_{i'i}\, (B_k)_{j'j}
    \]
    by the action property and orthonormality. This is the block-diagonal matrix with
    blocks $\Id_{m_k} \otimes B_k$.
\end{proof}

\begin{theorem}[Global unitary block-diagonal form]%
    \label{thm:starSubalgebra_unitary_block_diagonal}
    \lean{StarSubalgebra.exists_unitary_conj_blockDiagonal}
    \leanok
    Let $S$ be a $*$-subalgebra of $\MN{D}$. There are $K \in \N$, positive dimensions
    $d_0, \ldots, d_{K-1}$ and multiplicities $m_0, \ldots, m_{K-1}$ with
    $\sum_{k} d_k m_k = D$, realized by an explicit identification of the index sets, and
    a unitary $U \in \MN{D}$ such that every $A \in S$ satisfies
    \[
        U^{\dagger} A\, U \;=\; \bigoplus_{k=0}^{K-1} \Id_{m_k} \otimes B_k
    \]
    for matrices $B_k \in \MN{d_k}$ depending on $A$. Up to reordering the two tensor
    factors of each block this is the containment direction of the unital case of the
    block representation of~\cite[Theorem~6.14]{Wolf2012Quantum}: $S$ is carried by $U$
    into the block algebra $\bigoplus_k \Id_{m_k} \otimes \MN{d_k}$. The equality,
    including the reverse inclusion, is
    Theorem~\ref{thm:starSubalgebra_unitary_block_diagonal_iff}.
\end{theorem}

\begin{proof}
    \leanok
    \uses{thm:starSubalgebra_global_adapted_orthonormal_basis,
        lem:starSubalgebra_unitary_change_of_basis_block}
    Let $(f_{k,i,j})$ be the adapted orthonormal basis of
    Theorem~\ref{thm:starSubalgebra_global_adapted_orthonormal_basis} and let $U$ be the
    unitary supplied by Lemma~\ref{lem:starSubalgebra_unitary_change_of_basis_block} for
    this basis. Every $A \in S$ acts on the basis by matrices $B_k$ on the irreducible
    index, identically across the multiplicity index, so the lemma gives
    $U^{\dagger} A\, U = \bigoplus_k \Id_{m_k} \otimes B_k$.
\end{proof}

The reverse inclusion of the block representation rests on the double-commutant
characterization: in finite dimensions a $*$-subalgebra contains every matrix that commutes
with its commutant. The proof is the Jacobson density theorem for $\C^{D}$ as a semisimple
module over the subalgebra.

\begin{lemma}[Membership through the double commutant]%
    \label{lem:starSubalgebra_double_commutant}
    \lean{StarSubalgebra.mem_of_forall_comm}
    \leanok
    Let $S$ be a $*$-subalgebra of $\MN{D}$ and let $T \in \MN{D}$. If $T$ commutes with
    every linear operator on $\C^{D}$ that commutes with all members of $S$, then
    $T \in S$.
\end{lemma}

\begin{proof}
    \leanok
    \uses{lem:starSubalgebra_isSemisimpleModule}
    The space $\C^{D}$ is a semisimple module over $S$
    (Lemma~\ref{lem:starSubalgebra_isSemisimpleModule}), and the linear operators commuting
    with all members of $S$ are exactly the endomorphisms of this module. The hypothesis
    therefore makes $T$ linear over the endomorphism ring of the module. By the Jacobson
    density theorem, for every finite set of vectors there is a member of $S$ acting on
    them as $T$ does. Applied to a basis of $\C^{D}$, this produces a member of $S$ whose
    action agrees with $T$ on a basis, hence equals $T$.
\end{proof}

\begin{theorem}[Middle-factor coefficients of commuting overlapping operators]%
    \label{thm:commuting_overlapping_middle_coefficients}
    \lean{Matrix.middleSlices_commute_of_overlappingLifts_commute}
    \lean{Matrix.leftOverlappingLift}
    \lean{Matrix.rightOverlappingLift}
    \lean{Matrix.leftMiddleSlice}
    \lean{Matrix.rightMiddleSlice}
    \leanok
    Let $X_{AB}\in\mathbf L(H_A\otimes H_B)$ and
    $Y_{BC}\in\mathbf L(H_B\otimes H_C)$. For fixed outer indices define their
    middle-factor coefficient matrices by
    \[
        (X_{aa'})_{bb'}=(X_{AB})_{(a,b),(a',b')},\qquad
        (Y_{cc'})_{bb'}=(Y_{BC})_{(b,c),(b',c')}.
    \]
    If
    \[
        (X_{AB}\otimes\Id_C)(\Id_A\otimes Y_{BC})
        =(\Id_A\otimes Y_{BC})(X_{AB}\otimes\Id_C),
    \]
    then $X_{aa'}Y_{cc'}=Y_{cc'}X_{aa'}$ for every $a,a',c,c'$. This is the
    coefficientwise commutation step in the proof of
    \cite[Lemma~2.1]{Beigi2012Classification}.
\end{theorem}

\begin{proof}
    \leanok
    Evaluate the commutation identity between the basis vectors
    $\lvert a,b,c\rangle$ and $\lvert a',b',c'\rangle$. The identity matrices
    collapse the sums over the first and third intermediate indices, leaving
    \[
        \sum_{r}(X_{aa'})_{br}(Y_{cc'})_{rb'}
        =\sum_{r}(Y_{cc'})_{br}(X_{aa'})_{rb'}.
    \]
    This is precisely the $(b,b')$ entry of
    $X_{aa'}Y_{cc'}=Y_{cc'}X_{aa'}$.
\end{proof}

\begin{theorem}[Complementary spatial actions of a $*$-subalgebra and its commutant]%
    \label{thm:starSubalgebra_complementary_spatial_actions}
    \lean{StarSubalgebra.exists_complementary_action_orthonormalBasis}
    \leanok
    Let $S$ be a $*$-subalgebra of $\MN{D}$. There are positive integers
    $d_k,m_k$ and an orthonormal basis
    $(f_{k,i,j})_{k<K,\,i<m_k,\,j<d_k}$ of $\C^D$ such that every $A\in S$ acts by
    \begin{equation}\label{eq:spec_subalgebra_action}
        A f_{k,i,j}=\sum_{j'=0}^{d_k-1}(B_k)_{j'j}f_{k,i,j'},
    \end{equation}
    while every $T\in\MN{D}$ commuting with all members of $S$ acts by
    \begin{equation}\label{eq:spec_commutant_action}
        T f_{k,i,j}=\sum_{i'=0}^{m_k-1}(C_k)_{i'i}f_{k,i',j}.
    \end{equation}
    Thus the same orthonormal identification realizes $S$ on the second tensor factor
    and its commutant on the complementary first tensor factor. This is the
    finite-dimensional $C^*$-algebra step used in the proof of
    \cite[Lemma~2.1]{Beigi2012Classification}.
\end{theorem}

\begin{proof}
    \leanok
    \uses{thm:starSubalgebra_global_adapted_orthonormal_basis}
    Choose the adapted orthonormal basis $(f_{k,i,j})$ for $S$.
    Equation~(\ref{eq:spec_subalgebra_action}) is its defining action property.
    If $T$ commutes with $S$, the row transport
    \[
        \tau_{k,i',i}(v)
        =\sum_j\langle f_{k,i',j},v\rangle f_{k,i,j}
    \]
    also commutes with $S$, and so does $\tau_{k,i',i}T$. Schur's lemma makes this
    composite scalar on each irreducible row and makes its matrix coefficients vanish
    between different isotypic components. Expanding $T f_{k,i,j}$ in the orthonormal
    basis therefore gives equation~(\ref{eq:spec_commutant_action}), with
    coefficients
    $(C_k)_{i'i}$ independent of $j$.
\end{proof}

\begin{theorem}[Coefficient form of the commuting-overlap spatial decomposition]%
    \label{thm:commuting_overlapping_coefficient_spatial_decomposition}
    \lean{Matrix.exists_middle_spatial_decomposition_of_overlappingLifts_commute}
    \lean{Matrix.star_rightMiddleSlice}
    \leanok
    Let $X_{AB}\in\mathbf L(H_A\otimes H_B)$ and
    $Y_{BC}\in\mathbf L(H_B\otimes H_C)$, with $Y_{BC}$ Hermitian, and suppose
    \[
        [X_{AB}\otimes\Id_C,\Id_A\otimes Y_{BC}]=0.
    \]
    There is an orthonormal decomposition
    \[
        H_B\cong\bigoplus_{q=0}^{K-1} H_{q,r}\otimes H_{q,l}
    \]
    such that all middle-factor coefficients $X_{aa'}$ act only on $H_{q,l}$,
    while all middle-factor coefficients $Y_{cc'}$ act only on $H_{q,r}$. More
    precisely, in an orthonormal basis $(e_{q,r,s})$ adapted to this decomposition,
    \[
        X_{aa'}e_{q,r,s}
        =\sum_{s'}(B^{aa'}_q)_{s's}e_{q,r,s'},\qquad
        Y_{cc'}e_{q,r,s}
        =\sum_{r'}(C^{cc'}_q)_{r'r}e_{q,r',s}.
    \]
    The source lemma assumes that both overlapping operators are Hermitian; the
    coefficient conclusion requires Hermiticity only for $Y_{BC}$. This is the
    middle-space coefficient form of
    \cite[Lemma~2.1]{Beigi2012Classification}.
\end{theorem}

\begin{proof}
    \leanok
    \uses{thm:commuting_overlapping_middle_coefficients,
        thm:starSubalgebra_complementary_spatial_actions}
    Let $S$ be the commutant of the family $(Y_{cc'})_{c,c'}$ together with its
    adjoints. Hermiticity gives $Y_{cc'}^*=Y_{c'c}$.
    Theorem~\ref{thm:commuting_overlapping_middle_coefficients} therefore places
    every $X_{aa'}$ in $S$. Conversely, every $Y_{cc'}$ commutes with all members
    of $S$. Apply
    Theorem~\ref{thm:starSubalgebra_complementary_spatial_actions} to $S$: its
    members act on one tensor factor of each summand, and matrices commuting with
    $S$ act on the complementary factor. Specializing these two actions to
    $X_{aa'}$ and $Y_{cc'}$ gives the displayed formulas.
\end{proof}

\begin{theorem}[Explicit unitary block actions for commuting overlaps]%
    \label{thm:commuting_overlapping_unitary_block_actions}
    \lean{Matrix.exists_unitary_blockActions_of_overlappingLifts_commute}
    \lean{Matrix.leftSpatialBlockEquiv}
    \lean{Matrix.rightSpatialBlockEquiv}
    \leanok
    Let $X_{AB}\in\mathbf L(H_A\otimes H_B)$ and
    $Y_{BC}\in\mathbf L(H_B\otimes H_C)$ be Hermitian, and suppose
    \[
        [X_{AB}\otimes\Id_C,\Id_A\otimes Y_{BC}]=0.
    \]
    There are positive integers $d_q,m_q$, a unitary identification
    \[
        H_B\cong\bigoplus_{q=0}^{K-1}H_{q,l}\otimes H_{q,r},
    \]
    and Hermitian operators
    $R_q\in\mathbf L(H_A\otimes H_{q,l})$ and
    $S_q\in\mathbf L(H_{q,r}\otimes H_C)$ such that
    \[
        (\Id_A\otimes U^*)X_{AB}(\Id_A\otimes U)
        =\bigoplus_{q=0}^{K-1}(R_q\otimes\Id_{H_{q,r}}),
    \]
    and
    \[
        (U^*\otimes\Id_C)Y_{BC}(U\otimes\Id_C)
        =\bigoplus_{q=0}^{K-1}(\Id_{H_{q,l}}\otimes S_q).
    \]
    This is the explicit coordinate form of the two block-action identities in
    \cite[Lemma~2.1]{Beigi2012Classification}.
\end{theorem}

\begin{proof}
    \leanok
    \uses{thm:commuting_overlapping_coefficient_spatial_decomposition}
    Use the orthonormal decomposition from
    Theorem~\ref{thm:commuting_overlapping_coefficient_spatial_decomposition}.
    Let $U$ be the unitary whose columns are the adapted basis vectors.  Assemble the
    coefficient matrices $(B_q^{aa'})_{a,a'}$ into a matrix $R_q$ on
    $H_A\otimes H_{q,l}$, and similarly assemble $(C_q^{cc'})_{c,c'}$ into a matrix
    $S_q$ on $H_{q,r}\otimes H_C$.  Taking matrix entries in the adapted basis turns
    the coefficient-action formulas into the two displayed direct sums.  Finally,
    unitary conjugation preserves Hermiticity.  Restricting the two conjugated
    Hermitian operators to a summand, and fixing one basis vector in the nonempty
    complementary factor, shows that every $R_q$ and every $S_q$ is Hermitian.
\end{proof}

\begin{theorem}[Block-sum equalities for commuting overlapping operators]%
    \label{thm:commuting_overlapping_block_sum_equalities}
    \lean{Matrix.exists_unitary_blockSum_equalities_of_overlappingLifts_commute}
    \leanok
    \uses{thm:commuting_overlapping_unitary_block_actions}
    Let $X_{AB}\in\mathbf L(H_A\otimes H_B)$ and
    $Y_{BC}\in\mathbf L(H_B\otimes H_C)$ be Hermitian, and suppose
    \[
        [X_{AB}\otimes\Id_C,\Id_A\otimes Y_{BC}]=0.
    \]
    There are positive integers $d_q,m_q$, an identification
    \[
        H_B\cong\bigoplus_{q=0}^{K-1}H_{q,l}\otimes H_{q,r},
    \]
    a unitary $U\in\mathbf L(H_B)$, and Hermitian operators
    $R_q\in\mathbf L(H_A\otimes H_{q,l})$ and
    $S_q\in\mathbf L(H_{q,r}\otimes H_C)$ such that
    \[
        X_{AB}=(\Id_A\otimes U)
        \left(\bigoplus_{q=0}^{K-1}R_q\otimes\Id_{H_{q,r}}\right)
        (\Id_A\otimes U^*),
    \]
    and
    \[
        Y_{BC}=(U\otimes\Id_C)
        \left(\bigoplus_{q=0}^{K-1}\Id_{H_{q,l}}\otimes S_q\right)
        (U^*\otimes\Id_C).
    \]
    In these formulas the two direct sums are transported to the original
    product coordinates through the displayed identification.
    These are the two equalities in
    \cite[Lemma~2.1]{Beigi2012Classification}. The operators $R_q$ and $S_q$
    are Hermitian; they are not asserted to be unitary.
\end{theorem}

\begin{proof}
    \leanok
    Let $V_A=\Id_A\otimes U$ and $V_C=U\otimes\Id_C$. The preceding theorem
    gives the two conjugated operators in the direct-sum coordinates. Applying
    the inverse coordinate identifications yields
    \[
        V_A^*X_{AB}V_A
        =\bigoplus_q R_q\otimes\Id_{H_{q,r}},
        \qquad
        V_C^*Y_{BC}V_C
        =\bigoplus_q\Id_{H_{q,l}}\otimes S_q.
    \]
    Since $V_AV_A^*=\Id$ and $V_CV_C^*=\Id$,
    \[
        X_{AB}=V_A(V_A^*X_{AB}V_A)V_A^*,
        \qquad
        Y_{BC}=V_C(V_C^*Y_{BC}V_C)V_C^*,
    \]
    which are the stated equalities.
\end{proof}

Combining the containment direction with the double-commutant criterion yields the full
block representation of a finite-dimensional $*$-subalgebra, the unital case of
Eq.~(1.39) of~\cite{Wolf2012Quantum} invoked by~\cite[Theorem~6.14]{Wolf2012Quantum}: up
to a unitary change of basis, $S$ equals the block algebra
$\bigoplus_k \Id_{m_k} \otimes \MN{d_k}$.

\begin{theorem}[Full block form of a $*$-subalgebra]%
    \label{thm:starSubalgebra_unitary_block_diagonal_iff}
    \lean{StarSubalgebra.exists_unitary_conj_blockDiagonal_iff}
    \leanok
    Let $S$ be a $*$-subalgebra of $\MN{D}$. There are $K \in \N$, positive dimensions
    $d_0, \ldots, d_{K-1}$ and multiplicities $m_0, \ldots, m_{K-1}$ with
    $\sum_{k} d_k m_k = D$, realized by an explicit identification of the index sets, and
    a unitary $U \in \MN{D}$ such that a matrix $A \in \MN{D}$ belongs to $S$ exactly when
    \[
        U^{\dagger} A\, U \;=\; \bigoplus_{k=0}^{K-1} \Id_{m_k} \otimes B_k
    \]
    for some matrices $B_k \in \MN{d_k}$; that is, up to reordering the two tensor factors
    of each block,
    \[
        S \;=\; U \Bigl(\bigoplus_{k=0}^{K-1} \Id_{m_k} \otimes \MN{d_k}\Bigr) U^{\dagger}.
    \]
    This is the block representation of a finite-dimensional $*$-algebra in Eq.~(1.39)
    of~\cite{Wolf2012Quantum}, invoked by~\cite[Theorem~6.14]{Wolf2012Quantum}, in the
    unital case: a $*$-subalgebra contains the identity matrix, so there is no zero block.
\end{theorem}

\begin{proof}
    \leanok
    \uses{thm:starSubalgebra_global_adapted_orthonormal_basis,
        lem:starSubalgebra_unitary_change_of_basis_block,
        lem:starSubalgebra_double_commutant,
        lem:starSubalgebra_scalar_commutant_irreducible, def:starSubalgebra_intertwiner,
        def:starSubalgebra_sameIsotype}
    Let $(f_{k,i,j})$ be the adapted orthonormal basis of
    Theorem~\ref{thm:starSubalgebra_global_adapted_orthonormal_basis}, write $W_{k,i}$ for
    the span of the row $(f_{k,i,j})_{j < d_k}$, and let $U$ be the unitary supplied by
    Lemma~\ref{lem:starSubalgebra_unitary_change_of_basis_block} for this basis: any
    matrix $A$ whose action on the adapted basis is
    $A f_{k,i,j} = \sum_{j'} (B_k)_{j'j}\, f_{k,i,j'}$ satisfies the displayed conjugation
    identity; in particular every member of $S$ does, which is the containment direction.

    For the reverse inclusion, fix blocks $(B_k)_k$ and let $T$ be the matrix acting on
    the adapted basis by $T f_{k,i,j} = \sum_{j'} (B_k)_{j'j}\, f_{k,i,j'}$. By
    Lemma~\ref{lem:starSubalgebra_double_commutant} it suffices to show that $T$ commutes
    with every operator $g$ that commutes with all members of $S$; we compute the matrix
    of such a $g$ in the adapted basis. For a component index $k$ and multiplicity indices
    $i, i'$, consider the operator
    \[
        \tau \colon v \;\longmapsto\; \sum_{j} \langle f_{k,i',j}, v\rangle\, f_{k,i,j},
    \]
    which reads off the coefficients of a vector along the row $(k, i')$ and rebuilds them
    along the row $(k, i)$. Because every member of $S$ acts on the rows $(k, i')$ and
    $(k, i)$ by the same matrix, $\tau$ commutes with all members of $S$, and so does
    $\tau \circ g$. The composite maps $\C^{D}$ into $W_{k,i}$, and $W_{k,i}$ is
    irreducible, so by Schur's lemma
    (Lemma~\ref{lem:starSubalgebra_scalar_commutant_irreducible}) there is
    $\gamma^{(k)}_{i'i} \in \C$ with $\tau(g\,x) = \gamma^{(k)}_{i'i}\, x$ for all
    $x \in W_{k,i}$. Evaluating at $x = f_{k,i,j}$ and comparing coefficients gives
    \[
        \langle f_{k,i',j'},\, g\, f_{k,i,j}\rangle
        \;=\; \gamma^{(k)}_{i'i}\, \delta_{j'j}.
    \]
    Across components the matrix of $g$ vanishes: for $k' \neq k$,
    \[
        \langle f_{k',i',j'},\, g\, f_{k,i,j}\rangle \;=\; 0:
    \]
    with $\tau$ built from the rows $(k', i')$ of another component, the composite
    $\tau \circ g$ carries $W_{k,i}$ into $W_{k',i'}$ and commutes with $S$ on $W_{k,i}$,
    so a nonzero such entry would make it an intertwiner that does not vanish on
    $W_{k,i}$, giving $W_{k,i}$ and $W_{k',i'}$ the same type and contradicting
    Theorem~\ref{thm:starSubalgebra_global_adapted_orthonormal_basis}.
    Expanding $g f_{k,i,j}$ in the adapted basis therefore gives
    \[
        g\, f_{k,i,j} \;=\; \sum_{i'} \gamma^{(k)}_{i'i}\, f_{k,i',j} :
    \]
    in the adapted basis $g$ is block diagonal with blocks $C_k \otimes \Id_{d_k}$,
    opposite to the blocks of $S$. Since $T$ acts by one matrix $B_k$ on the irreducible
    index of every row and $g$ acts by one coefficient family $\gamma^{(k)}$ on the
    multiplicity index of every column, the two actions commute on every basis vector, so
    $T g = g T$. By Lemma~\ref{lem:starSubalgebra_double_commutant}, $T \in S$.

    Finally, if a matrix $A$ satisfies
    $U^{\dagger} A\, U = \bigoplus_k \Id_{m_k} \otimes B_k$, then
    $U^{\dagger} A\, U = U^{\dagger} T\, U$ for the member $T \in S$ built from the blocks
    $(B_k)_k$, so $A = T \in S$.
\end{proof}

\begin{theorem}[Fixed-point algebra is semisimple]%
    \label{thm:fixedPointAlgebra_isSemisimpleRing}
    \lean{Kraus.fixedPointAlgebra_isSemisimpleRing}
    \leanok
    \uses{thm:adjointFixedPointsStarSubalgebra}
    The adjoint-fixed-point $*$-subalgebra of a trace-preserving
    Kraus map on $\MN{D}$ is a semisimple ring.
\end{theorem}

\begin{proof}
    \leanok
    \uses{thm:adjointFixedPointsStarSubalgebra}
    The subalgebra is finite-dimensional (as a subalgebra of~$\MN{D}$),
    hence Artinian. Following~\cite[Theorem~6.14]{Wolf2012Quantum}, it
    suffices to show the Jacobson radical~$J$ vanishes. If $x \in J$, then
    $x^*x \in J$
    (left-ideal closure). The radical of an Artinian ring is nilpotent,
    so $x^*x$ is nilpotent. A self-adjoint nilpotent matrix satisfies
    $\lVert x^*x\rVert^{2^k} = \lVert (x^*x)^{2^k}\rVert = 0$,
    hence $x^*x = 0$, and the C$^*$-identity gives $x = 0$.
\end{proof}

\begin{theorem}[Abstract Wedderburn--Artin decomposition]%
    \label{thm:fixedPointAlgebra_wedderburnArtin}
    \lean{Kraus.fixedPointAlgebra_wedderburnArtin}
    \leanok
    \uses{thm:fixedPointAlgebra_isSemisimpleRing}
    There exist $n\in\N$ and dimensions $d_1,\ldots,d_n\ge 1$ such that
    the adjoint-fixed-point $*$-subalgebra is $\C$-algebra isomorphic to
    \[
        \prod_{k=1}^{n} \MN{d_k}.
    \]
\end{theorem}

\begin{proof}
    \leanok
    \uses{thm:fixedPointAlgebra_isSemisimpleRing}
    Combine semisimplicity (Theorem~\ref{thm:fixedPointAlgebra_isSemisimpleRing})
    with the Wedderburn--Artin structure theorem for finite-dimensional
    semisimple algebras over an algebraically closed field.
\end{proof}

\begin{theorem}[Wedderburn--Artin decomposition of the unital fixed-point algebra]%
    \label{thm:fixedPointsStarSubalgebra_exists_algEquiv_pi_matrix}
    \lean{Kraus.fixedPointsStarSubalgebra_exists_algEquiv_pi_matrix}
    \leanok
    \uses{thm:starSubalgebra_wedderburnArtin, thm:fixedPointsStarSubalgebra}
    Let $E$ be a unital Kraus map on $\MN{D}$ whose adjoint map $E^*$ has a
    positive definite fixed point $\rho > 0$. There exist $n \in \N$ and
    dimensions $d_1,\ldots,d_n \ge 1$ such that the fixed-point $*$-subalgebra
    $\mathrm{Fix}(E)$ is $\C$-algebra isomorphic to
    \[
        \mathrm{Fix}(E) \;\cong\; \prod_{k=1}^{n} \MN{d_k}.
    \]
\end{theorem}

\begin{proof}
    \leanok
    \uses{thm:starSubalgebra_wedderburnArtin, thm:fixedPointsStarSubalgebra}
    By Theorem~\ref{thm:fixedPointsStarSubalgebra} the fixed-point set
    $\mathrm{Fix}(E)$ is a $*$-subalgebra of $\MN{D}$. Applying the abstract
    Wedderburn--Artin decomposition for $*$-subalgebras
    (Theorem~\ref{thm:starSubalgebra_wedderburnArtin}) gives
    \[
        \mathrm{Fix}(E) \;\cong\; \prod_{k=1}^{n} \MN{d_k}
    \]
    with each block size $d_k \ge 1$.
\end{proof}

\begin{theorem}[Fixed-point algebra decomposition with multiplicities]%
    \label{thm:adjointFixedPoints_wedderburnDecomp}
    \lean{Kraus.adjointFixedPoints_wedderburnDecomp}
    \leanok
    \uses{thm:fixedPointAlgebra_wedderburnArtin}
    There exist integers $n$, block sizes $d_1,\ldots,d_n \ge 1$,
    multiplicities $m_1,\ldots,m_n \ge 1$, and a $\C$-algebra isomorphism from
    the adjoint-fixed-point $*$-subalgebra to
    \[
        \prod_{k=1}^{n} \MN{d_k}
    \]
    such that
    \[
        \sum_{k=1}^{n} d_k m_k \le D.
    \]
\end{theorem}
\begin{proof}
    \leanok
    \uses{thm:fixedPointAlgebra_wedderburnArtin}
    Apply the abstract Wedderburn decomposition
    (Theorem~\ref{thm:fixedPointAlgebra_wedderburnArtin}).
    Since the adjoint-fixed-point algebra is realized as a subalgebra of the
    ambient matrix algebra $\MN{D}$, this realization also yields
    multiplicities $m_k$ and the bound $\sum_k d_k m_k \le D$.
\end{proof}

\begin{remark}[Block form of the fixed-point algebra]
    Wolf's block form identifies the adjoint-fixed-point algebra, up to a
    unitary change of basis, with
    \[
        0 \oplus \bigoplus_k \MN{d_k} \otimes \Id_{m_k}
    \]
    together with the density-block form $\MN{d_k} \otimes \rho_k$.
    The theorem above gives the product decomposition, multiplicities, and
    dimension bound from this block form; the unitary conjugation identity itself, in
    the unital case, is Theorem~\ref{thm:adjointFixedPoints_block_form} below.
\end{remark}

\begin{theorem}[Unitary block form of the Heisenberg-picture fixed points]%
    \label{thm:adjointFixedPoints_block_form}
    \lean{Kraus.adjointFixedPoints_blockDiagonal_iff}
    \leanok
    \uses{def:adjoint_fixed_points, def:kraus_map_channels}
    Let $E$ be a trace-preserving Kraus map on $\MN{D}$ with a positive definite fixed
    point $\rho > 0$, $E(\rho) = \rho$. There are $n \in \N$, positive dimensions
    $d_0, \ldots, d_{n-1}$ and multiplicities $m_0, \ldots, m_{n-1}$ with
    $\sum_k d_k m_k = D$, realized by an explicit identification of the index sets, and
    a unitary $U \in \MN{D}$ such that a matrix $X \in \MN{D}$ satisfies $E^*(X) = X$
    exactly when
    \[
        U^{\dagger} X\, U \;=\; \bigoplus_{k=0}^{n-1} \Id_{m_k} \otimes B_k
    \]
    for some matrices $B_k \in \MN{d_k}$; that is, up to reordering the two tensor
    factors of each block,
    \[
        \mathrm{Fix}(E^*) \;=\;
        U \Bigl(\bigoplus_{k=0}^{n-1} \Id_{m_k} \otimes \MN{d_k}\Bigr) U^{\dagger}.
    \]
    The positive definite fixed point removes the zero block: this is the unital case
    of the block representation in Equation~(1.39) of~\cite{Wolf2012Quantum}, invoked
    by~\cite[Theorem~6.14]{Wolf2012Quantum}. The general form
    of~\cite[Theorem~6.14]{Wolf2012Quantum}, with a zero block and density weights
    $\rho_k$ on the Schr\"odinger-picture fixed points, is not asserted here.
\end{theorem}

\begin{proof}
    \leanok
    \uses{thm:adjointFixedPointsStarSubalgebra,
        thm:starSubalgebra_unitary_block_diagonal_iff}
    The fixed points of $E^*$ form a $*$-subalgebra of $\MN{D}$ because $E$ is
    trace-preserving with a positive definite fixed point
    (Theorem~\ref{thm:adjointFixedPointsStarSubalgebra}). The full block form of a
    $*$-subalgebra (Theorem~\ref{thm:starSubalgebra_unitary_block_diagonal_iff})
    supplies the unitary $U$, the dimensions and multiplicities, and the equivalence
    between membership and the block-diagonal conjugation identity.
\end{proof}

\begin{theorem}[Unitary block form of the fixed points of a unital Kraus map]%
    \label{thm:fixedPoints_block_form}
    \lean{Kraus.fixedPoints_blockDiagonal_iff}
    \leanok
    \uses{def:kraus_fixed_points, def:kraus_map_channels}
    Let $E$ be a unital Kraus map on $\MN{D}$ whose adjoint map $E^*$ has a positive
    definite fixed point $\rho > 0$. There are $n \in \N$, positive dimensions $d_k$
    and multiplicities $m_k$ with $\sum_k d_k m_k = D$, and a unitary $U \in \MN{D}$
    such that a matrix $X \in \MN{D}$ satisfies $E(X) = X$ exactly when
    \[
        U^{\dagger} X\, U \;=\; \bigoplus_{k=0}^{n-1} \Id_{m_k} \otimes B_k
    \]
    for some matrices $B_k \in \MN{d_k}$. This is
    Theorem~\ref{thm:adjointFixedPoints_block_form} with the roles of the map and its
    adjoint exchanged.
\end{theorem}

\begin{proof}
    \leanok
    \uses{thm:fixedPointsStarSubalgebra,
        thm:starSubalgebra_unitary_block_diagonal_iff}
    The fixed points of the unital map $E$ form a $*$-subalgebra of $\MN{D}$
    (Theorem~\ref{thm:fixedPointsStarSubalgebra}), and the full block form of a
    $*$-subalgebra (Theorem~\ref{thm:starSubalgebra_unitary_block_diagonal_iff})
    supplies the unitary and the equivalence.
\end{proof}

\begin{definition}[Canonical full-matrix extension of a direct-sum map]%
    \label{def:direct_sum_full_matrix_extension}
    \lean{Matrix.directSumDiagonalEmbedding}
    \lean{Matrix.directSumDiagonalCompression}
    \lean{Matrix.directSumExtension}
    \lean{Matrix.directSumMapExtension}
    \lean{Matrix.directSumTraceAdjointMap}
    \lean{Matrix.directSumTraceAdjointMapBetween}
    \leanok
    Let
    \[
        \mathcal A=\bigoplus_{k\in I}M_{n_k}(\mathbb C),
        \qquad
        \mathcal B=\bigoplus_{l\in J}M_{m_l}(\mathbb C),
    \]
    and let $\mathcal H_A=\bigoplus_{k\in I}\mathbb C^{n_k}$ and
    $\mathcal H_B=\bigoplus_{l\in J}\mathbb C^{m_l}$.  Write $\iota_A,\iota_B$
    for the block-diagonal embeddings and $\pi_A,\pi_B$ for diagonal-block
    compression.  For a linear map $T:\mathcal A\to\mathcal B$, its canonical
    full-matrix extension is
    \[
        \widehat T=\iota_B\circ T\circ\pi_A:
        \End(\mathcal H_A)\longrightarrow
        \End(\mathcal H_B).
    \]
    For $X\in\mathcal A$ and $Y\in\mathcal B$, the trace adjoint
    $T^*:\mathcal B\to\mathcal A$ is defined by
    \[
      \sum_{k\in I}\tr(T^*(Y)_kX_k)
      =\sum_{l\in J}\tr(Y_lT(X)_l).
    \]
    This is the block-diagonal realization needed for the fixed-point and
    classification arguments in
    \cite[Appendix~C.4, lines~1980--2003]{Cirac2016MPDO_arXiv}.
    It is only an auxiliary construction: it does not assert the
    classification conclusion of that passage.
\end{definition}

The preceding definition and the trace-adjoint laws in the next lemma provide
supporting facts for the classification argument.  By themselves, they do not
prove Schwarz equality or multiplicativity.  The inverse-CPTP theorem below
supplies those two conclusions; relabeling and dimension matching of the simple
summands, and implementation on each summand by unitary conjugation, remain
outside the present result.

\begin{lemma}[Trace-adjoint laws for two finite sums of matrix algebras]%
    \label{lem:direct_sum_rectangular_trace_adjoint}
    \lean{Matrix.sum_trace_directSumTraceAdjointMapBetween_mul}
    \lean{Matrix.directSumTraceAdjointMapBetween_involutive}
    \lean{Matrix.directSumTraceAdjointMapBetween_id}
    \lean{Matrix.directSumTraceAdjointMapBetween_comp}
    \lean{Matrix.directSumTraceAdjointMapBetween_self}
    \lean{Matrix.IsTracePreservingBetweenDirectSums.directSumTraceAdjointMapBetween_one}
    \lean{Matrix.traceAdjointMap_directSumMapExtension}
    \lean{IsKrausCP.traceAdjointMap}
    \lean{Matrix.IsKrausDirectSumMap.map_posSemidef}
    \lean{Matrix.IsKrausDirectSumMap.directSumTraceAdjointMapBetween}
    \leanok
    \uses{def:direct_sum_full_matrix_extension}
    For maps $T:\mathcal A\to\mathcal B$ and
    $S:\mathcal B\to\mathcal C$, trace adjoints satisfy
    \[
        \id_{\mathcal A}^*=\id_{\mathcal A},
        \qquad
        (S\circ T)^*=T^*\circ S^*,
        \qquad (T^*)^*=T,
        \qquad
        (\widehat T)^*=\widehat{T^*}.
    \]
    If $T$ preserves the total trace, then $T^*(\Id_{\mathcal B})=
    \Id_{\mathcal A}$.  If $\widehat T$ admits a Kraus representation, then
    so does $\widehat{T^*}$; in particular, both $T$ and $T^*$ send families
    of positive-semidefinite matrices to positive-semidefinite families.
    These statements are only the adjoint and positivity part of the
    classification argument; none of its multiplicative or blockwise
    conclusions is asserted here.
\end{lemma}

\begin{proof}
    \leanok
    \uses{def:direct_sum_full_matrix_extension}
    The defining trace identity and nondegeneracy of the trace pairing give
    the identity and composition formulas.  Total-trace preservation gives,
    for every $X\in\mathcal A$,
    \[
        \sum_{k\in I}\tr(T^*(\Id_{\mathcal B})_kX_k)
        =\sum_{l\in J}\tr(T(X)_l)
        =\sum_{k\in I}\tr(X_k),
    \]
    so $T^*(\Id_{\mathcal B})=\Id_{\mathcal A}$.  Finally, if
    $\widehat T(X)=\sum_a K_aXK_a^*$, cyclicity of the trace shows that
    \[
        (\widehat T)^*(Y)=\sum_a K_a^*YK_a.
    \]
    Thus the adjoint again has a Kraus representation, and compression of
    positive block-diagonal matrices gives the stated positivity.
\end{proof}

\begin{lemma}[Trace adjoint of a trace-preserving star isomorphism]
    \label{lem:direct_sum_trace_adjoint_star_equiv_inverse}
    \lean{Matrix.directSumTraceAdjointMapBetween_starAlgEquiv_eq_symm}
    \leanok
    Let $\Phi:\mathcal A\simeq\mathcal B$ be a star-algebra isomorphism
    between finite products of full matrix algebras.  If $\Phi$ preserves the
    total block trace, then
    \[
        \Phi^*=\Phi^{-1}.
    \]
\end{lemma}

\begin{proof}\leanok
    \uses{lem:direct_sum_rectangular_trace_adjoint}
    For $A\in\mathcal B$ and $B\in\mathcal A$, multiplicativity and trace
    preservation give
    \[
      \sum_j\tr(A_j\Phi(B)_j)
      =\sum_j\tr
        \bigl(\Phi(\Phi^{-1}(A)B)_j\bigr)
      =\sum_i\tr(\Phi^{-1}(A)_iB_i).
    \]
    Nondegeneracy of the trace pairing gives the asserted identity.
\end{proof}

\begin{theorem}[Channel and fixed-point compatibility of the canonical extension]%
    \label{thm:direct_sum_full_matrix_extension_compatibility}
    \lean{Matrix.IsPositiveDirectSumMap.directSumExtension_isPositiveMap}
    \lean{Matrix.IsSchwarzDirectSumMap.directSumExtension_isSchwarzMap}
    \lean{Matrix.IsTracePreservingDirectSumMap.directSumExtension_isTracePreservingMap}
    \lean{Matrix.traceAdjointMap_directSumExtension}
    \lean{Matrix.IsSchwarzDirectSumMap.traceAdjointMap_directSumExtension_isSchwarzMap}
    \lean{Matrix.directSumExtension_apply_eq_self_iff}
    \leanok
    \uses{def:direct_sum_full_matrix_extension}
    In the endomorphic case $\mathcal B=\mathcal A$, write
    $\mathcal H=\mathcal H_A=\mathcal H_B$ and
    $\iota=\iota_A=\iota_B$.
    If $T$ is positive and satisfies the Schwarz inequality on
    $\mathcal A$, then $\widehat T$ is positive and satisfies the Schwarz
    inequality on $\End(\mathcal H)$.  If $T$ preserves the
    total trace $\sum_k\tr(A_k)$, then $\widehat T$ preserves
    the ordinary trace.  Moreover,
    \[
        (\widehat T)^*=\widehat{T^*}.
    \]
    If $T$ is positive and its direct-sum trace adjoint satisfies the Schwarz
    inequality, then the trace adjoint of $\widehat T$ satisfies the Schwarz
    inequality on $\End(\mathcal H)$.  Finally,
    \[
        \Fix(\widehat T)=\iota(\Fix(T)).
    \]
    Equivalently, for every $X\in\End(\mathcal H)$,
    \[
        \widehat T(X)=X
        \quad\Longleftrightarrow\quad
        \exists A\in\mathcal A,\quad T(A)=A\ \text{ and }\ X=\iota(A).
    \]
\end{theorem}

\begin{proof}
    \leanok
    \uses{def:direct_sum_full_matrix_extension}
    For the matrix $R_k$ formed from the column blocks of $A$ outside the
    $k$th diagonal block, diagonal compression satisfies
    \[
        \pi(A^*A)_k-\pi(A)_k^*\pi(A)_k=R_k^*R_k\succeq0.
    \]
    Block-diagonal embedding is a positive algebra homomorphism and converts
    the sum of the block traces into the full trace.  It follows that
    positivity, both stated Schwarz implications, and trace preservation pass
    to $\widehat T$.  The embedding and compression satisfy
    \[
      \tr(\iota(A)X)
      =\sum_k\tr(A_k\pi(X)_k).
    \]
    Substitution of $\widehat T=\iota\circ T\circ\pi$ and the defining
    pairing for $T^*$ gives
    \[
        (\widehat T)^*=\widehat{T^*}.
    \]
    Since $\pi\circ\iota=\id$, the two fixed-point
    implications are
    \[
      \widehat T(X)=X
      \ \Longrightarrow\
      T(\pi(X))=\pi(\widehat T(X))=\pi(X),
      \qquad
      X=\widehat T(X)=\iota(T(\pi(X)))=\iota(\pi(X))
    \]
    and
    \[
      T(A)=A
      \ \Longrightarrow\
      \widehat T(\iota(A))
      =\iota(T(\pi(\iota(A))))=\iota(T(A))=\iota(A).
    \]
\end{proof}
